\documentclass[a4paper,12pt]{extarticle}

% Package imports
\usepackage{latexsym}
\usepackage{xcolor}
\usepackage{float}
\usepackage{ragged2e}
\usepackage[empty]{fullpage}
\usepackage{wrapfig}
\usepackage{lipsum}
\usepackage{tabularx}
\usepackage{titlesec}
\usepackage{geometry}
\usepackage{marvosym}
\usepackage{verbatim}
\usepackage{enumitem}
\usepackage{fancyhdr}
\usepackage{multicol}
\usepackage{graphicx}
\usepackage{cfr-lm}
\usepackage[T1]{fontenc}
\usepackage{fontawesome5}

% Color definitions
\definecolor{darkblue}{RGB}{0,0,139}
\definecolor{linkedinblue}{HTML}{0A66C2}
\definecolor{black}{RGB}{0,0,0}

% Page layout
\setlength{\multicolsep}{0pt} 
\pagestyle{fancy}
\fancyhf{} % clear all header and footer fields
\fancyfoot{}
\renewcommand{\headrulewidth}{0pt}
\renewcommand{\footrulewidth}{0pt}
\geometry{left=1.35cm, top=1.8cm, right=1.35cm, bottom=0.8cm}
\setlength{\footskip}{6pt} % Addressing fancyhdr warning

% Hyperlink setup (moved after fancyhdr to address warning)
\usepackage[hidelinks]{hyperref}
\hypersetup{
  colorlinks=true,
  linkcolor=black,
  filecolor=black,
  urlcolor=black,
}

% Custom box settings
\usepackage[most]{tcolorbox}
\tcbset{
    frame code={},
    center title,
    left=0pt,
    right=0pt,
    top=0pt,
    bottom=0pt,
    colback=gray!20,
    colframe=white,
    width=\dimexpr\textwidth\relax,
    enlarge left by=-2mm,
    boxsep=6pt,
    arc=0pt,outer arc=0pt,
}

% URL style
\urlstyle{same}

% Text alignment
\raggedright
\setlength{\tabcolsep}{0in}

% Section formatting
\titleformat{\section}{
  \scshape\raggedright\large
}{}{0em}{}[\color{black}\titlerule ]

% Custom commands
\newcommand{\resumeItem}[2]{
  \item{
    \textbf{#1}{\hspace{0.5mm}#2}
  }
}

\newcommand{\resumeSubheading}[4]{
  \item
    \begin{tabular*}{0.98\textwidth}[t]{l@{\extracolsep{\fill}}r}
        \textbf{#1} & \textit{\footnotesize{#4}} \\
        \textit{\footnotesize{#3}} &  \footnotesize{#2}\\
    \end{tabular*}
}

\newcommand{\resumeProject}[4]{
  \item
    \begin{tabular*}{0.98\textwidth}[t]{l@{\extracolsep{\fill}}r}
        \textbf{#1} & \textit{\footnotesize{#3}} \\
        \footnotesize{\textit{#2}} & \footnotesize{#4}
    \end{tabular*}

}

\newcommand{\resumeSubItem}[2]{\resumeItem{#1}{#2}}

\renewcommand{\labelitemi}{$\vcenter{\hbox{\tiny$\bullet$}}$}
\renewcommand{\labelitemii}{$\vcenter{\hbox{\tiny$\circ$}}$}

% Subheading list without bullets
\newcommand{\resumeSubHeadingListStart}{\begin{itemize}[leftmargin=0.3cm,label={},labelsep=0mm]}
\newcommand{\resumeHeadingSkillStart}{\begin{itemize}[leftmargin=0.3cm,itemsep=1.7mm, rightmargin=2ex]}
\newcommand{\resumeItemListStart}{\begin{itemize}[leftmargin=0.3cm,labelsep=1mm,itemsep=0.75mm]}

\newcommand{\resumeSubHeadingListEnd}{\end{itemize}}
\newcommand{\resumeHeadingSkillEnd}{\end{itemize}}
\newcommand{\resumeItemListEnd}{\end{itemize}}
\newcommand{\cvsection}[1]{%

\begin{tcolorbox}
    \textbf{\large #1}
\end{tcolorbox}
}

\newcolumntype{L}{>{\raggedright\arraybackslash}X}%
\newcolumntype{R}{>{\raggedleft\arraybackslash}X}%
\newcolumntype{C}{>{\centering\arraybackslash}X}%

% Commands for icon sizing and positioning
\newcommand{\socialicon}[1]{\raisebox{-0.05em}{\resizebox{!}{1em}{#1}}}
\newcommand{\mailicon}[1]{\raisebox{-0.08em}{\resizebox{!}{0.72em}{#1}}}
\newcommand{\ieeeicon}[1]{\raisebox{-0.3em}{\resizebox{!}{1.3em}{#1}}}

% Font options
\newcommand{\headerfonti}{\fontfamily{phv}\selectfont} % Helvetica-like (similar to Arial/Calibri)
\newcommand{\headerfontii}{\fontfamily{ptm}\selectfont} % Times-like (similar to Times New Roman)
\newcommand{\headerfontiii}{\fontfamily{ppl}\selectfont} % Palatino (elegant serif)
\newcommand{\headerfontiv}{\fontfamily{pbk}\selectfont} % Bookman (readable serif)
\newcommand{\headerfontv}{\fontfamily{pag}\selectfont} % Avant Garde-like (similar to Trebuchet MS)
\newcommand{\headerfontvi}{\fontfamily{cmss}\selectfont} % Computer Modern Sans Serif
\newcommand{\headerfontvii}{\fontfamily{qhv}\selectfont} % Quasi-Helvetica (another Arial/Calibri alternative)
\newcommand{\headerfontviii}{\fontfamily{qpl}\selectfont} % Quasi-Palatino (another elegant serif option)
\newcommand{\headerfontix}{\fontfamily{qtm}\selectfont} % Quasi-Times (another Times New Roman alternative)
\newcommand{\headerfontx}{\fontfamily{bch}\selectfont} % Charter (clean serif font)

\begin{document}
\headerfontii

% Header
\begin{center}
    {\fontsize{23pt}{18pt}\selectfont \textbf{Sunjae Lee}} \\[0.6em]
    \href{mailto:sl6100@columbia.edu}{\mailicon{\faEnvelope}\ sl6100@columbia.edu} |
    \href{mailto:satelite2517@gmail.com}{satelite2517@gmail.com} |
    \href{https://www.linkedin.com/in/sunjae-lee-81002727a/}{\textcolor{linkedinblue}{\raisebox{-0.08em}{\resizebox{!}{0.85em}{\faLinkedin}}}\ sunjae-lee}
    %\href{https://satelite2517.github.io/}{satelite2517.github.io}
\end{center}


\section{\textbf{Research Interests}}
\small{
  \begingroup{}
  \leftskip=0.3cm
  My research focuses on contributing to the realization of a \textcolor{darkblue}{\textbf{magnetic-confinement fusion reactor.}}
  Especially, I'm interested in these specific topics : 1) \textcolor{darkblue}{\textbf{kinetic MHD}} in 2) \textcolor{darkblue}{\textbf{negative triangularity(NT)}}, and 3) development of \textcolor{darkblue}{\textbf{AI-based high-performance simulations}}.
  % \textbf{Topics :} \textit{MHD Stability; Negative triangularity; Artificial Intelligence (AI); High-performance computing;}
  \par
  \endgroup
}

% \section{\textbf{Research Interests}}
% \small{
%   \begingroup{}
%   \leftskip=0.3cm
%   My research focuses on advancing the understanding of \textcolor{darkblue}{\textbf{magnetic-confinement fusion physics}} toward more efficient and stable plasma operation.
%   In particular, I am interested in (1) \textcolor{darkblue}{\textbf{MHD stability in negative-triangularity plasmas}} and (2) \textcolor{darkblue}{\textbf{AI-based high-performance plasma simulations}}.
%   \par
%   \endgroup
% }


\section{\textbf{Education}}
\resumeSubHeadingListStart

  \resumeSubheading
  {Seoul National University}{Seoul, Korea}
  {College of Liberal Studies / Dept. of Nuclear Engineering}{Mar 2021 – Jul 2026}
  \resumeItemListStart
    \item B.S. in \textbf{Computer Science \& Engineering} and \textbf{Physics} \hfill \textit{Feb 2026}
    \item M.S. in \textbf{Nuclear Engineering}, enrolled Mar 2026 and withdrew in Jul 2026
    \resumeSubItem{Relevant Coursework:}
    {Fusion Plasma Theory (Grad); Intro to Nuclear Fusion; E\&M; Thermal \& Stat. Physics; Mathematical Methods of Physics;
    Scalable HPC (Grad); Numerical Analysis; Intro to ML; Deep Learning}
  \resumeItemListEnd

  \resumeSubheading
  {Columbia University in the City of New York}{New York, NY, USA}
  {Applied Physics and Applied Mathematics (APAM)}{Beginning Fall 2026}
  \resumeItemListStart
    \item M.S./Ph.D. in \textbf{Applied Physics (Plasma Physics)}
  \resumeItemListEnd

\resumeSubHeadingListEnd

\section{\textbf{Research Experiences}}
  \resumeSubHeadingListStart
    \resumeSubheading
    {Dept. of Nuclear Engineering, Seoul National University}{Seoul, Korea}
    {Undergraduate Research Intern, 3DST Lab}{Sep 2024 – Present}
    % {{}[\href{https://volunteer-org-a-link.com}{\textcolor{darkblue}{\faIcon{globe}}}]}
    \resumeItemListStart
      \item[] \href{https://nucleng.snu.ac.kr/about/people/faculty?mode=view&profidx=167}{\textbf{3 Dimensional Stellarator \& Tokamak Lab - Advisor: Prof. Jong-Kyu Park}}
      \item Conducted stability analysis of NT using ideal, drift-kinetic, and full kinetic MHD models with GPEC suite
      \item Found NT has stronger kinetic stabilization effect due to closed-orbit trapped particle's resonance 
      \item Identified a second-stability region of the $n = 1$ in DIII-D NT and analyzed the effect of magnetic shear
    \resumeItemListEnd

    \resumeSubheading
      {Princeton Plasma Physics Laboratory (PPPL)}{Princeton, NJ, USA}
      {Research Student, SNU–PPPL Summer Collaboration}{Jul 2025 – Aug 2025}
      % {{}[\href{https://volunteer-org-a-link.com}{\textcolor{darkblue}{\faIcon{globe}}}]}
    \resumeItemListStart
      \item Participated in the PPPL–SNU Summer Research Collaboration, focusing on nonlinear 3D MHD modeling, transport analysis, and pilot-plant design for NSTX-U and KSTAR plasmas
      \item Learned and applied the BALOO code for edge-stability analysis in the K-DEMO project and NT 
    \resumeItemListEnd

    \resumeSubheading
    {Dept. of Nuclear Engineering, Seoul National University}{Seoul, Korea}
    {Undergraduate Research Intern, NUPLEX Lab}{Sep 2023 – Aug 2024}
    \resumeItemListStart
      \item[] \href{https://nucleng.snu.ac.kr/about/people/faculty?mode=view&profidx=169}{\textbf{Seoul National University Plasma \& Ion Beam Laboratory - Advisor: Prof. Y. S. Hwang}}
      \item Designed and implemented a database for the Versatile Experiment Spherical Torus(VEST) using the IMAS–HSDS data architecture
      \item Developed a post-processing tool (VAFT) for automated analysis of experimental diagnostic shot data
      \item Explored operational space for the first time, identifying performance limits from current-driven instabilities
    \resumeItemListEnd

    % \resumeSubheading
    %   {Dept. of Energy Resoureces Engineering, Seoul National University }{Seoul, Korea}
    %   {Undergraduate Research Participant}{Apr 2022 – Dec 2022}
    %   \resumeItemListStart
    %     \item[] \href{https://eee.snu.ac.kr/}{Electrochemical Energy Engineering Laboratory - Advisor: Prof. Jin soo Kang}
    %     \item Conducted CDI-based lithium recovery experiments and assessed scalability via efficiency comparison
    % \resumeItemListEnd 
    
  \resumeSubHeadingListEnd

%%%%%

\section*{\textbf{Publications} \hfill {\scriptsize \textit{J: Journal, C: Conference}}}
\small{\begin{enumerate}[leftmargin=*, labelsep=0.5em, align=left, widest={[\textbf{S.1}]}, itemindent=0.3cm, label={\textbf{[\arabic*]}}]

\item[\textbf{[J.1]}] Hong-Sik Yun$^{\dagger}$, \textbf{Sunjae Lee$^{\dagger}$}, Laurent Jung, et al. (2025). \href{https://doi.org/10.1088/1361-6587/ae1b6a}
{\textbf{Developing an IMAS-Compatible Platform for the University-Level Tokamak VEST and Its Application in Operating Characteristics Analysis}}. 
\textit{Plasma Physics and Controlled Fusion}, Vol. 67, DOI:10.1088/1361-6587/ae1b6a
%  Vol. 00, Issue 0, pp. 000-000.

$^{\dagger}$ These authors contributed equally to this work.
\end{enumerate}
}

\section*{\textbf{Orals \& Posters}\hfill {\scriptsize \textit{P: Poster presentations, O: Oral presentations}}}
\small{
  
\begin{enumerate}[leftmargin=*, labelsep=0.5em, align=left, widest={[\textbf{S.1}]}, itemindent=0em, label={\textbf{[\arabic*]}}]

\item[\textbf{[P.1]}] \textbf{S. J. Lee}, J. B. Cho, O. Nelson, C. Paz-Soldan, J.-K. Park. (2025). \href{https://satelite2517.github.io/files/2025_Nov_APS-DPP.pdf}{\textbf{Kinetic MHD Stability of Low-n Modes in Negative Triangularity Plasmas}}.  
\textit{67th Annual Meeting of the APS Division of Plasma Physics}, Long Beach, CA, USA (Nov 2025) 
% \href{https://schedule.aps.org/dpp/2025/events/BP13/138}{\textcolor{darkblue}{(Abstract)}}

\item[\textbf{[P.2]}] \textbf{S. J. Lee}, J. B. Cho, O. Nelson, C. Paz-Soldan, J.-K. Park. (2025).\href{https://satelite2517.github.io/files/2025_Jul_MHD_workshop.pdf}{\textbf{Kinetic stability of negative-triangularity plasmas}}. 
\textit{29th Workshop on MHD Stability Control}, Princeton, New Jersey, USA (Jul 2025)

\item[\textbf{[P.3]}] \textbf{S. J. Lee}, J.-K. Park, H.S. Yun, L. Jung, G. W. Nam, G. G. Seo, J. M. Lee (2025). \href{https://satelite2517.github.io/files/2025_Apr_KPS.pdf}{\textbf{Standardized Data Infrastructure for Tokamak: Implementation in VEST (Versatile Experiment spherical tokamak)}}. 
\textit{2025 Spring meeting of the Korean Physics Society}, Seoul, Korea (Apr 2025)

\item[\textbf{[P.4]}] \textbf{S. J. Lee}, H.S. Yun, Laurent Jung, Jung-Hwa Kim, Y.S. Hwang (2024). \href{https://satelite2517.github.io/files/2024_Jun_IFPC.pdf}{\textbf{Improvements of database system and analysis suite in VEST}}. 
\textit{3rd International Fusion and Plasma conference (IFPC)}, Seoul, Korea (Jun 2024)

\end{enumerate}
}

% \section{\textbf{Awards \& Honors}}
% \small{
%   \begingroup
%   \leftskip=0.3cm
%   \textbf{Home Renovation Volunteer}, Habitat for Humanity Korea (\href{https://www.instagram.com/snuhabitat.official/}{Snuhabitat}) \hfill Sep 2021 – Feb 2022\\
%   \par
%   \endgroup
% }

\section{\textbf{Projects}}
  \resumeSubHeadingListStart

  \resumeProject
    {\href{https://github.com/OpenFUSIONToolkit/JPEC}{Perturbed Equilibrium Code Hackathon – Columbia Fusion Research Center}}
    {Tools: Julia, Fortran, GPEC (reference), Plasma Physics Modeling}
    {Jul 2025}
    {Columbia, NY, USA [\href{https://fusion.columbia.edu/news/inaugural-hackathon-hosted-fusion-research-center}{\textcolor{darkblue}{\faIcon{globe}}}]}
  \resumeItemListStart
    \item Participated in a hackathon to reproduce DCON’s energy-principle-based ideal MHD stability solver in Julia
    \item Developed coordinate transformation and numerical routines to compute stability analysis from equilibrium
  \resumeItemListEnd

  \resumeProject
    {\href{https://github.com/satelite2517/Accelerator_summer_school}{High-Performance CUDA Programming Project – Accelerator Programming Summer School}}
    {Tools: CUDA, NVIDIA GPU, Nsight, C++}
    {Aug 2024}
    {Seoul, Korea}
  \resumeItemListStart
    \item Implemented parallel programs in CUDA/C++ for high-throughput computation in GPU 
    \item Analyzed performance using Nsight Systems, identifying bottlenecks in memory access and execution flow
  \resumeItemListEnd

  \resumeProject
    {Text-to-Image Generation with Conditional GAN - 2023 OUTTA Deep Learning Bootcamp}
    {Tools: PyTorch, Conditional GAN}
    {Jun 2023 – Aug 2023}
    {Seoul, Korea}
  \resumeItemListStart
    \item Built a text-to-image generation model using conditional GAN architecture and Ranked 2 out of 60 teams
  \resumeItemListEnd

  \resumeSubHeadingListEnd


\section{\textbf{Teaching Experience}}
\small{
  \begingroup
  \leftskip=0.3cm
  \textbf{Teaching Assistant}, VEST Winter Camp — SNU \hfill Feb 3–11, 2026\\
  \hspace*{0.3cm}\footnotesize $\circ$ Tokamak stability course for undergraduate and graduate students\\[0.25em]
  \textbf{Teaching Assistant}, Computer Architecture — SNU \hfill Fall 2025\\
  \hspace*{0.3cm}\footnotesize $\circ$ Undergraduate-level course
  \par
  \endgroup
}

\section{\textbf{Leadership \& Volunteer Experiences}}
\small{
  \begingroup
  \leftskip=0.3cm
  \textbf{Vice President}, \href{https://www.instagram.com/snu_sorijigi/}{Sorijigi (Classical Music Listening Club)} – SNU \hfill Sep 2021 – Feb 2023\\
  \textbf{Volunteer}, \href{https://www.instagram.com/snuhabitat.official/}{Habitat for Humanity Korea (Snuhabitat) – SNU} \hfill Sep 2021 – Feb 2022\\
  \textbf{Volunteer}, \href{https://www.instagram.com/tail_snu/}{Tail – SNU} \hfill Aug 2022 – Dec 2023\\
  \par
  \endgroup
}


\section{\textbf{Skills \& Additional Informations}}
 \resumeHeadingSkillStart
  \resumeSubItem{Programming Languages:}
    {Python, C/C++, Fortran, Julia, Linux, CUDA, MATLAB, OpenCL, MPI, OpenMP}
  \resumeSubItem{Simulation Tools:}
    {DCON, GPEC, BALOO, KineticEFITtime}
  \resumeSubItem{TOEFL:}
    {100/120 (Reading 27, Listening 28, Speaking 24, Writing 21)}
 \resumeHeadingSkillEnd

\end{document}
